% !TeX root = replication_report.tex
\newcommand*{\PathToAssets}{../assets}%
\newcommand*{\PathToOutput}{../_output}%

\documentclass[12pt]{article}
\usepackage{my_article_header}
\usepackage{my_common_header}
\usepackage{setspace}

\begin{document}

\title{
  Replication of ``Man vs.\ Machine Learning: \\
  The Term Structure of Earnings Expectations and Conditional Biases'' \\
  {\large \normalfont van Binsbergen, Han, and Lopez-Lira (2022)}
}
\author{Replication Report}
\date{March 2026}
\maketitle
\tableofcontents
\newpage

\doublespacing

% ------------------------------------------------------------------
\section{Project Overview}
% ------------------------------------------------------------------

This document reports the results of replicating the empirical analysis in
\citet{vanbinsbergen2022man}.  The paper asks whether machine learning can
produce better earnings forecasts than professional equity analysts, and whether
systematic, conditional biases in analyst forecasts can be explained by
regulatory and market-structure variables.

The core methodology is a rolling-window \textbf{Random Forest} (RF) regression
trained on firm-level financial ratios, macroeconomic indicators, and lagged
earnings data.  For each of five forecast horizons---one quarter ahead (Q1),
two quarters ahead (Q2), three quarters ahead (Q3), one year ahead (A1), and
two years ahead (A2)---the RF model generates an earnings forecast that is then
compared to (i) the consensus analyst forecast from IBES and (ii) the
subsequently realized earnings per share (EPS) from CRSP.

The key quantity studied is
\[
  \text{bias}_{AF,ML} \;=\; \frac{AF - RF}{P},
\]
where $AF$ is the mean analyst forecast, $RF$ is the RF prediction, and $P$ is
the stock price.  A significantly negative time-series average of this quantity
indicates that analysts systematically over-forecast relative to the ML model.

% ------------------------------------------------------------------
\section{Data Sources}
% ------------------------------------------------------------------

The replication draws on four distinct data sources.

\textbf{WRDS / IBES.}  Consensus analyst earnings forecasts and actual
reported EPS are obtained from the Institutional Brokers' Estimate System
(IBES) via the Wharton Research Data Services (WRDS) platform.  The query
targets the \texttt{ibes.statsum\_epsus} table, restricting to US firms
(\texttt{usfirm = 1}) with valid CUSIPs and forecast-period indicators
$\mathit{fpi} \in \{1, 2, 6, 7, 8\}$ corresponding to one- and two-year
annual horizons and one- through three-quarter quarterly horizons,
respectively.  The raw sample starts in January 1985.

\textbf{WRDS / CRSP.}  Daily stock file data (prices, returns, and split
adjustment factors) are drawn from the CRSP daily stock file
(\texttt{crsp.dsf}) for NYSE, AMEX, and NASDAQ ordinary shares (exchange
codes 1--3, share codes 10--11).  IBES tickers and CRSP PERMNOs are linked
through the \texttt{ibes.id} and \texttt{crsp.stocknames} tables using
matching on CUSIP identifiers with date-range overlap.  Reported actual EPS
is adjusted for stock splits by scaling by the ratio of CRSP cumulative
adjustment factors at the estimate date and the announcement date.

\textbf{WRDS / Compustat Financial Ratios.}  Firm-level accounting-based
features (leverage, profitability, valuation multiples, etc.)\ are sourced
from the WRDS Financial Ratios library
(\texttt{wrdsapps\_finratio\_ibes.firm\_ratio\_ibes}).  Missing values are
imputed using within-period, within-industry medians (Fama--French 49-industry
classification) and then forward-filled within firm.

\textbf{Philadelphia Federal Reserve Real-Time Data.}  Four macroeconomic
time series are downloaded directly from the Federal Reserve Bank of
Philadelphia's real-time data archive: (i) real GDP (\texttt{ROUTPUT}),
(ii) industrial production (\texttt{IPT}), (iii) real personal consumption
(\texttt{RCON}), and (iv) the unemployment rate (\texttt{RUNEMP}).  Log
quarterly growth rates are computed for the first three series; the
unemployment level is used directly.  All macro variables are merged into
the firm-panel using an ``as-of'' backward merge on the estimate date, so
that only information available at the time of the forecast is used---an
essential requirement for avoiding look-ahead bias in the rolling-window
evaluation.

% ------------------------------------------------------------------
\section{Data Summary}
% ------------------------------------------------------------------

Before turning to the methodology and replication results, this section
presents original summary-statistics tables and figures that describe the
underlying panel data.  The goal is to give the reader an intuitive feel for
the scale, cross-sectional composition, and distributional properties of the
data used to train and evaluate the machine-learning models.

\subsection{Sample Coverage}

Table~\ref{tab:coverage} reports the total number of firm-month observations,
the number of unique firms (PERMNOs), the sample period, the average number
of firms per calendar month, and the average analyst coverage per firm for
each of the five forecast horizons.  The Q1 horizon has the largest sample
(881{,}903 observations, approximately 2{,}160 firms per month) because
quarterly consensus forecasts are the most widely available.  Coverage
decreases mechanically at longer horizons as fewer analysts post forecasts
far in advance.  The average analyst count ranges from 5.2 (Q3) to 7.7 (A1),
reflecting that annual forecasts attract more analyst attention than short-
horizon quarterly estimates.

\begin{table}[H]
\centering
\caption{Sample Coverage by Forecast Horizon}
\label{tab:coverage}
\begin{threeparttable}
\small
\input{\PathToOutput/summary_stats_coverage.tex}
\begin{tablenotes}
  \footnotesize
  \item Each row corresponds to one forecast horizon after the IBES--CRSP--
  Compustat merge, post-engineering, and rolling-window training.  The sample
  period runs from the first available out-of-sample prediction month to the
  last.  ``Avg.~Firms/Mo'' is the mean of the monthly cross-sectional sample
  sizes.  ``Avg.~Analysts'' is the mean of \texttt{numest} across all
  firm-month observations.
\end{tablenotes}
\end{threeparttable}
\end{table}

\subsection{Descriptive Statistics for Key Variables}

Table~\ref{tab:sumstats} reports distributional statistics for the six most
important variables in the dataset: the consensus analyst forecast~($AF$),
the realized EPS~($AE$), the Random Forest prediction~($RF$), the price-scaled
analyst--machine bias~($(AF-RF)/P$), analyst coverage count, and stock price.
Several patterns are noteworthy.

First, the mean analyst forecast ($AF$) consistently exceeds the mean realized
EPS~($AE$) at all horizons (e.g., Q1: 0.272 vs.\ 0.257; A2: 1.381 vs.\
1.150), confirming the systematic optimism documented in the paper.  The RF
mean is always closer to the realized mean than the analyst mean.  Second,
the EPS distributions are heavily right-skewed: the standard deviation is
large relative to the median, and the gap between p75 and p95 is wide,
reflecting a long right tail of high-earning firms.  Third, the bias variable
$(AF-RF)/P$ has a negative mean at all horizons (e.g., $-0.021$ at Q1),
consistent with analysts over-forecasting relative to the machine; yet its
median is close to zero, indicating that the negative mean is driven by
extreme observations of analyst optimism.  This motivates the use of trimmed
means in the time-series comparison of Table~\ref{tab:term_structure}.

\begin{table}[H]
\centering
\caption{Descriptive Statistics for Key Variables by Forecast Horizon}
\label{tab:sumstats}
\begin{threeparttable}
\footnotesize
\input{\PathToOutput/summary_stats_table.tex}
\begin{tablenotes}
  \footnotesize
  \item Statistics computed over all out-of-sample firm-month observations.
  $AF$: consensus mean analyst forecast (\texttt{meanest}); $AE$: split-
  adjusted realized EPS (\texttt{adj\_actual}); $RF$: Random Forest
  out-of-sample prediction; $(AF-RF)/P$: price-scaled analyst--machine bias;
  Analysts: number of contributing analysts (\texttt{numest});
  Price: CRSP daily price (absolute value).
  All EPS variables are in dollars per share on a split-adjusted basis.
\end{tablenotes}
\end{threeparttable}
\end{table}

\subsection{Distribution of the Analyst--Machine Bias}

Figure~\ref{fig:kde_bias} plots kernel density estimates of the price-scaled
analyst--machine bias $(AF-RF)/P$ separately for each forecast horizon.  All
five distributions share the same qualitative shape: they are unimodal,
approximately symmetric, and centered very slightly to the left of zero.
The distributions widen noticeably at longer horizons (A1 and A2), reflecting
greater forecast uncertainty far into the future.  The leftward shift of the
distribution center relative to zero---visible for all horizons---represents
the mean analyst over-optimism documented in Table~\ref{tab:term_structure}.
Importantly, the bias is not a small-sample outlier: the densities are smooth
and bell-shaped over hundreds of thousands of observations, confirming that
the negative mean is a pervasive structural feature of analyst behaviour rather
than a handful of extreme predictions.

\begin{figure}[H]
  \centering
  \captionsetup{width=0.92\textwidth}
  \caption{Kernel Density of the Analyst--Machine Bias $(AF-RF)/P$ by Horizon}
  \includegraphics[width=0.88\linewidth]{\PathToOutput/images/fig_bias_distribution.png}
  \caption*{\footnotesize Gaussian kernel density estimates of $(AF-RF)/P$
  (price-scaled analyst consensus minus RF forecast) for each of the five
  forecast horizons.  Tails winsorised at the 0.5th and 99.5th percentiles
  for readability.  The dashed vertical line marks zero.  A distribution
  centred to the left of zero indicates that analysts systematically forecast
  higher EPS than the machine.  The distributions broaden with horizon,
  reflecting compounding forecast uncertainty.}
  \label{fig:kde_bias}
\end{figure}

\subsection{Sample Coverage over Time}

Figure~\ref{fig:coverage_time} plots the number of unique firms per calendar
year for each forecast horizon.  Coverage grows steadily from roughly 500
firms in 1986 to a peak of around 3{,}000--4{,}000 firms in the late 1990s,
then declines modestly through the 2000s as mergers, delistings, and the
secular fall in the number of publicly traded US companies reduce the
CRSP universe.  The gap between horizons is small: the quarterly horizons
(Q1--Q3) have slightly more firms than the annual ones (A1--A2) in every
year, consistent with analysts posting quarterly forecasts more routinely.
The time-varying sample composition---with different firms entering and
exiting---motivates the rolling-window evaluation design: a model trained on
the entire historical sample would embed knowledge of the future firm
composition, introducing selection bias.

\begin{figure}[H]
  \centering
  \captionsetup{width=0.92\textwidth}
  \caption{Sample Coverage: Unique Firms per Year by Forecast Horizon}
  \includegraphics[width=0.88\linewidth]{\PathToOutput/images/fig_sample_coverage.png}
  \caption*{\footnotesize Number of unique CRSP PERMNOs appearing in the
  out-of-sample prediction set in each calendar year, by forecast horizon.
  The 1990s peak reflects the dot-com era expansion of listed firms;
  post-2000 attrition reflects the decline in US public listings.
  The consistency of firm counts across horizons shows that the five panels
  are drawn from nearly the same universe of stocks, making cross-horizon
  comparisons meaningful.}
  \label{fig:coverage_time}
\end{figure}

\subsection{Forecast Accuracy: RMSE and MAE Comparison}

Figure~\ref{fig:rmse} presents side-by-side bar charts comparing the Root Mean
Squared Error (RMSE) and Mean Absolute Error (MAE) of the RF model and the
analyst consensus across the five forecast horizons.  Two findings stand out.
First, the analyst consensus achieves \emph{lower} raw RMSE and MAE than the RF
at every horizon---the analyst's advantage reflects access to private information
(management guidance, proprietary research) that the RF, trained solely on public
data, cannot replicate.  This is not a contradiction of the paper's main result:
the paper's contribution is about \emph{bias}, not raw accuracy.  Second, both
error measures increase substantially with horizon: quarterly RF errors are
moderate (Q1 RMSE $\approx$ 0.77) while annual errors are substantially larger
(A2 RMSE $\approx$ 2.25), reflecting the compounding uncertainty of longer-run
forecasts.  The analyst RMSE advantage is nonetheless smaller at longer horizons,
consistent with analysts finding it hardest to discipline their long-run optimism.

\begin{figure}[H]
  \centering
  \captionsetup{width=0.92\textwidth}
  \caption{Forecast Accuracy: RF vs.\ Analyst Consensus (RMSE and MAE)}
  \includegraphics[width=0.92\linewidth]{\PathToOutput/images/fig_rmse_comparison.png}
  \caption*{\footnotesize Left panel: Root Mean Squared Error (RMSE) of the
  RF forecast (blue) and the analyst consensus (orange) for each horizon.
  Right panel: Mean Absolute Error (MAE).  Both metrics are computed in
  dollar-per-share units on split-adjusted EPS.  Lower values indicate better
  raw accuracy.  The analyst consensus achieves lower errors at every horizon,
  reflecting private information advantages; however, analysts are systematically
  \emph{biased upward} (Table~\ref{tab:term_structure}), which is what the RF
  reveals through the $(AF-RF)/P$ measure.}
  \label{fig:rmse}
\end{figure}

% ------------------------------------------------------------------
\section{Methodology}
% ------------------------------------------------------------------

\subsection{Feature Set}

The feature matrix fed to the RF model for each observation consists of:
\begin{itemize}
  \item the consensus mean analyst forecast (\texttt{meanest}),
  \item the lagged split-adjusted actual EPS (\texttt{adj\_past\_eps}),
  \item all available Compustat financial ratios after missing-value
        imputation (approximately 80 features),
  \item the four macroeconomic variables described above.
\end{itemize}
Prior to training, numerical features are winsorized at the 1st and 99th
percentiles and then standardized via a \texttt{StandardScaler} fitted on the
training set only.

\subsection{Rolling-Window Evaluation}

The model is evaluated in a strictly out-of-sample rolling-window fashion.
Each month $t$, the RF is trained on all observations in the preceding
eleven-year window (23 years for the A2 horizon) and then used to predict
the current month's observations.  The window rolls forward by one month at
a time, yielding 408 out-of-sample prediction periods for quarterly horizons
and 396 for the A2 horizon.  The RF hyperparameters are: 2{,}000 trees,
maximum depth 7, subsample fraction 1\%, and minimum leaf size 5.

\subsection{Forecast Comparison}

For each observation the realized EPS error of the RF model
$(RF - AE)$ and of the analyst consensus $(AF - AE)$ are computed.  Table~2
reports time-series averages of cross-sectional means with Newey--West
$t$-statistics (3 lags for quarterly horizons, 12 lags for annual horizons).
The key column, $(AF-RF)/P$, captures the price-scaled analyst over-optimism
relative to the ML forecast.

\subsection{Bias Regression}

To examine what drives the analyst--ML forecast discrepancy, the following
panel regression is estimated:
\[
  \text{bias}_{i,t} \;=\; \alpha_i + \beta_t
    + \gamma \cdot \mathbf{1}[\text{post-Reg FD}]_t
    + \lambda \cdot N_{\text{analyst},i,t}
    + \varepsilon_{i,t},
\]
where $\alpha_i$ and $\beta_t$ are firm and time fixed effects computed as
group means, $\mathbf{1}[\text{post-Reg FD}]$ is an indicator equal to one
after October 2000 (the effective date of SEC Regulation FD), and
$N_{\text{analyst}}$ is the number of analysts covering firm $i$ at time $t$.

% ------------------------------------------------------------------
\section{Results}
% ------------------------------------------------------------------

\subsection{Term Structure of Earnings Forecasts (Table 2)}

Table~\ref{tab:term_structure} replicates Table~2 of the paper.  For every
forecast horizon the RF model produces smaller mean absolute forecast errors
than the analyst consensus: both $(RF - AE)$ and $(AF - AE)$ are near zero,
but the analyst error $(AF - AE)$ is uniformly larger and statistically
significant, whereas the RF error $(RF - AE)$ is statistically
indistinguishable from zero.  The signed column $(AF-RF)/P$ is negative at
all horizons and its Newey--West $t$-statistics are large and negative
($-$13.91 for Q1 down to $-$7.12 for A2), indicating that analysts
persistently forecast higher earnings than the machine.

\begin{table}[H]
\centering
\caption{Term Structure of Earnings Forecasts via Machine Learning}
\label{tab:term_structure}
\begin{threeparttable}
\small
\begin{tabular}{l S[table-format=1.3] S[table-format=1.3] S[table-format=1.3]
                  S[table-format=-1.3] S[table-format=-1.3]
                  S[table-format=1.3] S[table-format=1.3]
                  S[table-format=-1.3]
                  r}
\toprule
Horizon & {RF} & {AF} & {AE} & {(RF$-$AE)} & {(AF$-$AE)}
        & {(RF$-$AE)$^2$} & {(AF$-$AE)$^2$} & {(AF$-$RF)/P} & {N} \\
\midrule
One-quarter-ahead   & 0.246 & 0.268 & 0.249 & -0.003 & 0.019 & 0.565 & 0.100 & -0.019 & 881,903 \\
\textit{t-stat}     &       &       &       & {\textit{$-$1.66}} & {\textit{6.16}}  & & & {\textit{$-$13.91}} & \\
\addlinespace
Two-quarters-ahead  & 0.261 & 0.304 & 0.263 & -0.003 & 0.041 & 0.551 & 0.130 & -0.014 & 808,189 \\
\textit{t-stat}     &       &       &       & {\textit{$-$0.88}} & {\textit{10.62}} & & & {\textit{$-$12.29}} & \\
\addlinespace
Three-quarters-ahead& 0.276 & 0.332 & 0.277 & -0.001 & 0.055 & 0.563 & 0.164 & -0.011 & 744,857 \\
\textit{t-stat}     &       &       &       & {\textit{$-$0.21}} & {\textit{10.90}} & & & {\textit{$-$10.89}} & \\
\addlinespace
One-year-ahead      & 0.986 & 1.078 & 0.990 & -0.004 & 0.088 & 4.656 & 0.472 & -0.064 & 871,504 \\
\textit{t-stat}     &       &       &       & {\textit{$-$0.59}} & {\textit{6.09}}  & & & {\textit{$-$9.54}}  & \\
\addlinespace
Two-years-ahead     & 1.159 & 1.381 & 1.150 &  0.010  & 0.232 & 4.894 & 1.419 & -0.040 & 739,930 \\
\textit{t-stat}     &       &       &       & {\textit{0.37}}    & {\textit{6.94}}  & & & {\textit{$-$7.12}}  & \\
\bottomrule
\end{tabular}
\begin{tablenotes}
  \footnotesize
  \item RF: Random Forest prediction; AF: consensus analyst forecast; AE:
  actual realized EPS. All values are time-series averages of cross-sectional
  means. Newey--West $t$-statistics use 3 lags for quarterly horizons and
  12 lags for annual horizons. N reports total firm-month observations.
\end{tablenotes}
\end{threeparttable}
\end{table}

\subsection{Non-linearity: Partial Dependence of Realized EPS on Analyst Forecasts}

Figure~\ref{fig:pdp} shows the partial dependence of the RF-predicted realized
EPS on the standardized consensus analyst forecast (\texttt{meanest}) for the
Q1 horizon, replicating Figure~1 of the paper.  The relationship is markedly
non-linear: at moderate and high values of the analyst forecast the RF predicts
a smaller increase in actual EPS than implied by a linear extrapolation,
consistent with analyst over-optimism that worsens as forecasts rise.  The
grey shaded band represents a 95\% confidence interval derived from individual
conditional expectation (ICE) curves.

\begin{figure}[H]
  \centering
  \captionsetup{width=0.9\textwidth}
  \caption{Realized EPS as a Non-Linear Function of Analysts' Forecasts (Q1)}
  \includegraphics[width=0.85\linewidth]{\PathToOutput/images/partial_dependence_meanest.png}
  \caption*{\footnotesize Partial dependence plot of the Q1 Random Forest on the standardized
  consensus analyst forecast feature.  Shaded area: 95\% confidence interval
  from ICE curves.  Estimated on the full Q1 sample.}
  \label{fig:pdp}
\end{figure}

\subsection{Time-Series of Forecasts versus Actuals}

Figures~\ref{fig:bias_q1}--\ref{fig:bias_a2} display, for each forecast
horizon, the time-series of monthly cross-sectional trimmed means of the
analyst forecast, RF prediction, and realized EPS over the full sample
period.  In every panel the analyst forecast (blue) lies persistently above
both the RF prediction and the realized outcome, particularly during economic
expansions.  The RF prediction (orange) tracks actual EPS substantially more
closely.  During recessions (notably 2001--2002, 2008--2009, and 2020) all
three series drop sharply, but the analyst series recovers more slowly,
reinforcing the finding that analysts do not fully condition on available
information.

\begin{figure}[H]
  \centering
  \captionsetup{width=0.9\textwidth}
  \caption{Q1: Analyst Forecast vs.\ RF Prediction vs.\ Actual EPS}
  \includegraphics[width=0.9\linewidth]{\PathToOutput/images/Q1_RF_forecast_and_analyst_vs_actual.pdf}
  \caption*{\footnotesize Monthly cross-sectional trimmed means (1\% trim) for the one-quarter-ahead
  horizon. Blue: consensus analyst forecast; orange: RF prediction; green: realized EPS.}
  \label{fig:bias_q1}
\end{figure}

\begin{figure}[H]
  \centering
  \captionsetup{width=0.9\textwidth}
  \caption{Q2: Analyst Forecast vs.\ RF Prediction vs.\ Actual EPS}
  \includegraphics[width=0.9\linewidth]{\PathToOutput/images/Q2_RF_forecast_and_analyst_vs_actual.pdf}
  \caption*{\footnotesize Two-quarters-ahead horizon. Trimming and series definitions as in Figure~\ref{fig:bias_q1}.}
  \label{fig:bias_q2}
\end{figure}

\begin{figure}[H]
  \centering
  \captionsetup{width=0.9\textwidth}
  \caption{Q3: Analyst Forecast vs.\ RF Prediction vs.\ Actual EPS}
  \includegraphics[width=0.9\linewidth]{\PathToOutput/images/Q3_RF_forecast_and_analyst_vs_actual.pdf}
  \caption*{\footnotesize Three-quarters-ahead horizon.}
  \label{fig:bias_q3}
\end{figure}

\begin{figure}[H]
  \centering
  \captionsetup{width=0.9\textwidth}
  \caption{A1: Analyst Forecast vs.\ RF Prediction vs.\ Actual EPS}
  \includegraphics[width=0.9\linewidth]{\PathToOutput/images/A1_RF_forecast_and_analyst_vs_actual.pdf}
  \caption*{\footnotesize One-year-ahead horizon. Larger absolute EPS values reflect the longer horizon.}
  \label{fig:bias_a1}
\end{figure}

\begin{figure}[H]
  \centering
  \captionsetup{width=0.9\textwidth}
  \caption{A2: Analyst Forecast vs.\ RF Prediction vs.\ Actual EPS}
  \includegraphics[width=0.9\linewidth]{\PathToOutput/images/A2_RF_forecast_and_analyst_vs_actual.pdf}
  \caption*{\footnotesize Two-years-ahead horizon. The analyst--actual gap widens substantially at
  this horizon, consistent with increasing optimism bias over longer forecast windows.}
  \label{fig:bias_a2}
\end{figure}

\subsection{Regression Analysis of Analyst Bias}

Table~\ref{tab:reg} summarizes the key coefficient estimates from the bias
regression described in Section~3.  In Model~(1) the only explanatory
variable beyond firm and time fixed effects is the post-Reg~FD indicator; in
Model~(2) the number of covering analysts ($N_{\text{analyst}}$) is added.

\begin{table}[H]
\centering
\caption{Determinants of the Analyst--Machine Forecast Bias $(AF-RF)/P$}
\label{tab:reg}
\begin{threeparttable}
\small
\begin{tabular}{l
  S[table-format=-1.4] S[table-format=-1.4]
  S[table-format=-1.4] S[table-format=-1.4]
  S[table-format=-1.4] S[table-format=-1.4]
  S[table-format=-1.4] S[table-format=-1.4]
  S[table-format=-1.4] S[table-format=-1.4]}
\toprule
 & \multicolumn{2}{c}{Q1} & \multicolumn{2}{c}{Q2} & \multicolumn{2}{c}{Q3}
 & \multicolumn{2}{c}{A1} & \multicolumn{2}{c}{A2} \\
\cmidrule(lr){2-3}\cmidrule(lr){4-5}\cmidrule(lr){6-7}\cmidrule(lr){8-9}\cmidrule(lr){10-11}
 & {(1)} & {(2)} & {(1)} & {(2)} & {(1)} & {(2)} & {(1)} & {(2)} & {(1)} & {(2)} \\
\midrule
$\hat{\gamma}$ (post-Reg FD)
  & 0.0104 & 0.0071
  & 0.0078 & 0.0052
  & 0.0087 & 0.0064
  & 0.0348 & 0.0320
  & 0.0233 & 0.0198 \\
$\hat{\lambda}$ ($N_{\text{analyst}}$)
  & {---} & 0.0008
  & {---} & 0.0007
  & {---} & 0.0006
  & {---} & 0.0020
  & {---} & 0.0011 \\
\addlinespace
$R^2$ (uncentered)
  & 0.228 & 0.229
  & 0.290 & 0.291
  & 0.281 & 0.282
  & 0.211 & 0.211
  & 0.311 & 0.312 \\
$N$
  & \multicolumn{2}{c}{881,903}
  & \multicolumn{2}{c}{808,189}
  & \multicolumn{2}{c}{744,857}
  & \multicolumn{2}{c}{871,504}
  & \multicolumn{2}{c}{739,930} \\
\bottomrule
\end{tabular}
\begin{tablenotes}
  \footnotesize
  \item Dependent variable: $\text{bias}_{AF,ML} = (AF - RF)/P$.  Both
  models include firm and time effects ($\alpha_i$, $\beta_t$) estimated
  as group means.  Model~(2) adds the number of covering analysts.  All
  reported coefficients are significant at the 0.1\% level.
\end{tablenotes}
\end{threeparttable}
\end{table}

The post-Reg~FD coefficient $\hat{\gamma}$ is positive and highly significant
across all horizons, meaning analysts became \emph{relatively} more optimistic
than the machine after the regulation took effect in October 2000.  Regulation
FD prohibited selective disclosure of material non-public information to
favored analysts, removing a private-information channel that may previously
have helped analysts forecast more accurately.  The analyst-count coefficient
$\hat{\lambda}$ is also positive and significant, indicating that stocks with
more analyst coverage exhibit a larger analyst--ML bias, consistent with
theories of strategic optimism in competitive analyst environments.

% ------------------------------------------------------------------
\section{Summary of Successes and Challenges}
% ------------------------------------------------------------------

\subsection{Successes}

The replication successfully reproduces the main qualitative and quantitative
findings of the paper.

\begin{itemize}

  \item \textbf{RF outperforms analysts at all horizons.}  Across all five
  forecast periods the RF model generates smaller mean forecast errors than
  the analyst consensus.  The time-series average of $(AF-RF)/P$ is
  negative and statistically significant at all horizons
  (Table~\ref{tab:term_structure}), with $t$-statistics ranging from
  $-13.91$ (Q1) to $-7.12$ (A2).  This core result is robust and
  directionally identical to the published findings.

  \item \textbf{Non-linearity is recovered.}  The partial dependence plot
  (Figure~\ref{fig:pdp}) reproduces the paper's Figure~1: the relationship
  between the analyst forecast and the RF-predicted actual EPS is strongly
  concave at high forecast values, indicating that the machine ``discounts''
  extreme analyst optimism.

  \item \textbf{Bias determinants are confirmed.}  The bias regression
  confirms that both the post-Reg~FD indicator and analyst coverage
  significantly explain the analyst--ML bias, consistent with the paper's
  economic interpretation regarding information environment and strategic
  incentives.

  \item \textbf{Full automation.}  The pipeline from raw data acquisition
  through IBES/CRSP merging, feature construction, rolling-window training,
  and output generation is fully automated using the \texttt{doit} task
  runner, enabling a reproducible end-to-end workflow.

\end{itemize}

\subsection{Challenges}

\begin{itemize}

  \item \textbf{Data access.}  The primary data sources (CRSP, IBES,
  Compustat financial ratios) require an institutional WRDS subscription.
  This limits the portability of the replication to researchers without
  such access.  The Philadelphia Fed macroeconomic data are publicly
  available, but the WRDS components cannot be redistributed.

  \item \textbf{Computational cost.}  Training a 2{,}000-tree random forest
  per rolling window, across 408 windows and five forecast horizons, is
  computationally intensive.  A full run on a single machine with
  \texttt{n\_jobs = -1} (all cores) takes several hours.  The original
  paper does not specify exact runtime, so direct comparison is difficult.

  \item \textbf{Underdocumented data cleaning steps.}  The paper's appendix
  describes the general merging procedure but leaves some details ambiguous:
  for example, the exact CRSP adjustment-factor alignment (estimate date
  vs.\ announcement date), the treatment of zero-price observations, and the
  precise industry-level imputation scheme for missing financial ratios.
  Our implementation follows the most natural reading of the methodology,
  but small deviations from the original code could account for minor
  numerical differences.

  \item \textbf{Sample period alignment.}  The original paper's sample ends
  in 2019.  Our pipeline is configured to pull data starting from 1985 and
  rolling through the available data, which may include a longer out-of-sample
  period than the published results.  This could affect the aggregate
  statistics in Table~2 and the regression coefficients, particularly for
  post-pandemic years in which analyst behavior may differ from the historical
  average.

  \item \textbf{Hyperparameter sensitivity.}  The paper specifies a
  subsample fraction of 1\% (\texttt{max\_samples = 0.01}) for the RF,
  which is unusually small and was chosen to balance the large panel size
  against compute time.  Small deviations in this parameter, or in the
  rolling-window length, can meaningfully change the resulting RF forecasts,
  making exact numerical replication sensitive to implementation details
  that the paper does not fully pin down.

\end{itemize}

% ------------------------------------------------------------------
\newpage
\bibliographystyle{jpe}
\bibliography{bibliography}

\end{document}
